% =============================================================================
% 06-process-states-schedulers-notes.tex — Lecture Notes for Lecture 6:
% Process State Diagram, Schedulers, and Context Switch
%
% DERIVED ARTIFACT. Source of truth: Slides/CS502/06-process-states-schedulers.tex.
% Generated by /lecture-notes; do not hand-edit content independently of the
% Beamer source without re-running /qa-notes afterward (single-source-of-truth.md).
% Structured as a textbook chapter (numbered sections/figures/examples,
% topic-organized rather than a frame-by-frame transcript). All TikZ figures
% are copied verbatim from the deck.
%
% Anchor reading: Silberschatz, Galvin & Gagne, Operating System Concepts,
%                 8th ed. — Ch. 3 §3.1 (process states, p. 103) and §3.2
%                 (process scheduling, pp. 105-110), page-verified via
%                 master_supporting_docs/CS502/supporting_books/SilberschatzOSC/index.md;
%                 Stallings, 6th ed. — Ch. 3 §3.2 (p. 111), §3.4 (p. 135);
%                 Bovet & Cesati, 3rd ed. — Ch. 3 (pp. 79-130).
%
% Compile with: cd Notes/CS502 &&
%   TEXINPUTS=../../Preambles:$TEXINPUTS xelatex -interaction=nonstopmode 06-process-states-schedulers-notes.tex  (x3)
% =============================================================================

\documentclass[11pt]{article}
\usepackage[margin=1in]{geometry}
\usepackage{array}
\input{../../Preambles/header}
\coursecode{CS502}

% Chapter-style numbering: this is Lecture 6's chapter.
\renewcommand{\thesection}{6.\arabic{section}}
\renewcommand{\thefigure}{6.\arabic{figure}}
\renewcommand{\thetable}{6.\arabic{table}}

\newtheorem{example}{Example}
\renewcommand{\theexample}{6.\arabic{example}}

\newenvironment{definitionbox}[1]{%
  \par\noindent\textbf{Definition (#1).}\ \itshape
}{\par\vspace{0.3em}}

\setlength{\emergencystretch}{2.5em}

\title{Process State Diagram, Schedulers, and Context Switch\\\large Lecture Notes --- Lecture 6\\
  \normalsize How the OS Tracks and Switches Processes}
\author{Auqib Hamid Lone\\[0.3em]
  \emph{Department of Computer Science \& Engineering}, \emph{GCET Ganderbal}}
\date{\today}

\begin{document}
\maketitle

\section{One CPU, Three Programs}

A lab machine has \emph{one} CPU core. On it, \texttt{gcc} is compiling a large
project, a text editor is waiting for the next keystroke, and a backup job has
been queued for the night. To the person at the keyboard, all three feel as if
they are running. The question that organizes this chapter is what that
feeling hides: \emph{at any instant, how many of them actually are running ---
and how does the operating system decide which one runs, remember where the
others stopped, and switch between them?}

We give the three programs names and keep them for the whole chapter:
\texttt{P1} is \texttt{gcc}, which has long CPU bursts; \texttt{P2} is the
editor, which has short bursts and spends most of its time waiting;
\texttt{P3} is the backup job, which has been submitted but not yet admitted.
On one core, at most \key{one} process is running at any instant. Everything
else in this chapter explains how the kernel manages the others.

\paragraph{What we already know.} Lecture~5 defined a \key{process} as a
program in execution: its code, data, stack, and heap, together with the CPU
state that goes with them. The kernel records each process in a \key{process
control block} (PCB), which holds the process state, the program counter
(\texttt{PC}), the registers, scheduling information, memory-management
information, accounting information, and the list of open files. The PCB is
what lets a process \emph{stop} and later \emph{resume} exactly where it left
off (Silberschatz et al., Ch.~3 \S3.1, pp.~101--103
\cite{SilberschatzOSC_operating_system_concepts}). The PCB has a \emph{state}
field. This chapter asks which values that field can hold, which events change
it, and who changes it.

\paragraph{Roadmap.} Four questions frame the chapter:
\begin{enumerate}
  \item \key{Which states can a process be in?} The five-state diagram and its
        legal transitions (Section~\ref{sec:states}).
  \item \key{Where does the kernel keep processes?} The job, ready, and device
        queues, and the queueing diagram (Section~\ref{sec:queues}).
  \item \key{Who moves processes between queues?} The long-term, short-term,
        and medium-term schedulers (Section~\ref{sec:schedulers}).
  \item \key{What does a switch cost?} The context switch, the dispatcher, and
        its overhead (Section~\ref{sec:cs}).
\end{enumerate}
The division of labour is worth stating once: the states are the vocabulary,
the queues are the data structures, the schedulers are the policies, and the
context switch is the mechanism.

\section{The Process State Diagram}
\label{sec:states}

\noindent\muted{\textit{Reading: Silberschatz et al., Ch.~3 \S3.1, p.~103;
Stallings, Ch.~3 \S3.2, p.~111.}}

A process does not simply ``run.'' It is created, waits its turn, runs, blocks,
and finishes, and the kernel tracks every one of those steps in a single field
of the PCB.

\begin{definitionbox}{Process state}
The state of a process is the PCB field that records its current activity. A
process is in exactly one of five states at any instant.
\end{definitionbox}

\begin{table}[h]
\centering
\small
\renewcommand{\arraystretch}{1.25}
\begin{tabular}{@{}>{\ttfamily}l>{\raggedright\arraybackslash}p{0.7\linewidth}@{}}
\toprule
\normalfont\textbf{State} & \textbf{Meaning} \\
\midrule
new & being created; the PCB is being set up \\
ready & in memory, able to run, waiting only for a CPU \\
running & its instructions are being executed on a CPU \\
waiting & blocked until some event occurs (I/O completion, a signal) \\
terminated & finished; its resources are being reclaimed \\
\bottomrule
\end{tabular}
\caption{The five process states (Silberschatz et al., Ch.~3 \S3.1, p.~103).}
\label{tab:states}
\end{table}

The names vary between books. Stallings calls \texttt{waiting} ``blocked'' and
\texttt{terminated} ``exit'' (Ch.~3 \S3.2, p.~111
\cite{StallingsOS_internals_design_principles}); the states themselves are the
same.

Figure~\ref{fig:five-state} draws the states as nodes and the permitted changes
as arrows. Read it from left to right: a process is admitted from
\texttt{new} into \texttt{ready}; the scheduler dispatches it to
\texttt{running}; from \texttt{running} it can be interrupted back to
\texttt{ready}, block in \texttt{waiting}, or exit to \texttt{terminated}; and
a waiting process returns to \texttt{ready} when its event completes.

\begin{figure}[h]
\centering
\begin{tikzpicture}[every node/.style={font=\small},
      st/.style={dag-node, minimum width=2.3cm, minimum height=0.8cm}]
    \node[st] (new) at (-5.6,1.4) {\texttt{new}};
    \node[st] (rdy) at (-2.2,0) {\texttt{ready}};
    \node[st] (run) at (2.2,0) {\texttt{running}};
    \node[st] (trm) at (5.6,1.4) {\texttt{terminated}};
    \node[st] (wt)  at (0,-2.1) {\texttt{waiting}};
    \draw[observed-edge] (new) -- node[above right, font=\scriptsize] {admitted} (rdy);
    \draw[observed-edge] (rdy.20) -- node[above, font=\scriptsize] {scheduler dispatch} (run.160);
    \draw[observed-edge] (run.200) -- node[below, font=\scriptsize] {interrupt} (rdy.340);
    \draw[observed-edge] (run) -- node[below right=-2pt, font=\scriptsize, align=left] {I/O or\\event wait} (wt);
    \draw[observed-edge] (wt) -- node[below left=-2pt, font=\scriptsize, align=right] {I/O or event\\completion} (rdy);
    \draw[observed-edge] (run) -- node[above left, font=\scriptsize] {exit} (trm);
  \end{tikzpicture}
\caption{The five-state process diagram --- after Silberschatz et al., Ch.~3
\S3.1, p.~103 \cite{SilberschatzOSC_operating_system_concepts}.}
\label{fig:five-state}
\end{figure}

\subsection{The six legal transitions}

Each arrow in Figure~\ref{fig:five-state} has a triggering event and an agent
that carries it out. Table~\ref{tab:transitions} lists all six.

\begin{table}[h]
\centering
\small
\renewcommand{\arraystretch}{1.25}
\begin{tabular}{@{}>{\raggedright\arraybackslash}p{0.27\linewidth}>{\raggedright\arraybackslash}p{0.32\linewidth}>{\raggedright\arraybackslash}p{0.32\linewidth}@{}}
\toprule
\textbf{Transition} & \textbf{Triggering event} & \textbf{Who acts} \\
\midrule
\texttt{new} $\to$ \texttt{ready} & admitted to the system & long-term scheduler \\
\texttt{ready} $\to$ \texttt{running} & chosen for the CPU & short-term scheduler + dispatcher \\
\texttt{running} $\to$ \texttt{ready} & timer interrupt / preemption & kernel (interrupt handler) \\
\texttt{running} $\to$ \texttt{waiting} & I/O request, \texttt{wait}, sleep & the process itself (system call) \\
\texttt{waiting} $\to$ \texttt{ready} & the awaited event occurs & kernel (I/O-completion interrupt) \\
\texttt{running} $\to$ \texttt{terminated} & \texttt{exit} or fatal error & the process / the kernel \\
\bottomrule
\end{tabular}
\caption{The six legal transitions, their triggers, and their agents.}
\label{tab:transitions}
\end{table}

Two transitions are conspicuously absent, and both absences are deliberate.
There is no \texttt{waiting} $\to$ \texttt{running} arrow: a process whose
event has just occurred must queue for the CPU like every other ready process.
There is no \texttt{ready} $\to$ \texttt{waiting} arrow: only a running process
can issue the request it would wait on, so a process that is not running has
nothing to block on.

The separation in the first case matters. Suppose \texttt{P2} is waiting for
the disk, the disk interrupt arrives, and the CPU happens to be idle. Does
\texttt{P2} go straight to \texttt{running}? No: in the diagram it goes
\texttt{waiting} $\to$ \texttt{ready}, because the interrupt handler only marks
it runnable. The short-term scheduler then chooses among \emph{all} ready
processes; with an idle CPU it may pick \texttt{P2} almost immediately, but it
is still a separate decision. Keeping ``wake up'' and ``get the CPU'' as two
steps is exactly what lets the scheduler enforce a policy.

\begin{example}[The editor, for one keystroke]
\label{ex:keystroke}
Follow \texttt{P2}, the editor, through a single keystroke.
\begin{enumerate}
  \item \texttt{P2} is \texttt{running}. It calls \texttt{read} on the
        keyboard, but no key is pending. Transition:
        \texttt{running} $\to$ \texttt{waiting}.
  \item The kernel places \texttt{P2}'s PCB on the keyboard's device queue and
        runs another process. \texttt{P2} now consumes no CPU time.
  \item You press a key. The keyboard interrupt handler finds \texttt{P2}
        waiting for it. Transition: \texttt{waiting} $\to$ \texttt{ready}.
  \item \texttt{P2}'s PCB joins the ready queue. It does \emph{not} run yet;
        whatever was running continues.
  \item The short-term scheduler picks \texttt{P2}, and the dispatcher loads
        its context. Transition: \texttt{ready} $\to$ \texttt{running}.
  \item \texttt{P2} echoes the character and calls \texttt{read} again.
        Transition: \texttt{running} $\to$ \texttt{waiting}.
\end{enumerate}
\emph{Result:} each keystroke costs the editor three transitions
(\texttt{waiting} $\to$ \texttt{ready} $\to$ \texttt{running} $\to$
\texttt{waiting}), and between keystrokes the editor spends almost all of its
life in \texttt{waiting}.
\end{example}

\subsection{How many processes can be in each state?}

With $N$ CPU cores, at most $N$ processes can be \texttt{running}. Any number
may be \texttt{ready} or \texttt{waiting} at once. On our one-core lab machine,
a snapshot at one instant might show \texttt{P1} \texttt{running},
\texttt{P2} \texttt{waiting} on the keyboard, and \texttt{P3} \texttt{new},
still in the batch queue.

\begin{example}[A snapshot of 40 processes on 4 cores]
A 4-core machine has 40 processes. At most 4 are \texttt{running}. If 9 are
\texttt{ready}, then the remaining $40 - 4 - 9 = 27$ are \texttt{waiting} (or
\texttt{new} or \texttt{terminated}). This is typical of desktops: most
processes are waiting most of the time.
\end{example}

\subsection{Suspended processes and the seven-state model}

Memory is finite. If every process in memory is \texttt{waiting}, the CPU idles
even though more work is queued on disk. The remedy is to \key{swap out} a
process: copy its image to disk and free its memory. Such a process is
\key{suspended}. Stallings (Ch.~3 \S3.2, p.~111) splits each of \texttt{ready}
and \texttt{waiting} into an in-memory version and a suspended version, giving
a \key{seven-state model} with two new states:

\begin{itemize}
  \item \texttt{ready/suspend}: the process is on disk but able to run once it
        is brought back into memory;
  \item \texttt{blocked/suspend}: the process is on disk \emph{and} still
        waiting for an event.
\end{itemize}

Figure~\ref{fig:seven-state} shows the model. The top row lives in memory; the
bottom row has been swapped out to disk. A blocked process can be suspended;
if its event occurs while it is on disk it moves to \texttt{ready/suspend}
without being brought back; and either suspended state can be activated back
into memory.

\begin{figure}[h]
\centering
\begin{tikzpicture}[every node/.style={font=\scriptsize},
      st/.style={dag-node, minimum width=2.2cm, minimum height=0.7cm, font=\scriptsize}]
    \node[st] (new) at (-5.7,1.6) {\texttt{new}};
    \node[st] (rdy) at (-1.8,1.6) {\texttt{ready}};
    \node[st] (run) at (2.4,1.6) {\texttt{running}};
    \node[st] (ext) at (5.7,1.6) {\texttt{exit}};
    \node[st] (blk) at (0.3,0) {\texttt{blocked}};
    \node[st] (rs)  at (-4.2,-1.7) {\texttt{ready/suspend}};
    \node[st] (bs)  at (2.4,-1.7) {\texttt{blocked/suspend}};
    \draw[observed-edge] (new) -- node[above] {admit} (rdy);
    \draw[observed-edge] (rdy.15) -- node[above] {dispatch} (run.165);
    \draw[observed-edge] (run.195) -- node[below, pos=0.3] {timeout} (rdy.345);
    \draw[observed-edge] (run) -- node[right, pos=0.6] {event wait} (blk);
    \draw[observed-edge] (blk) -- node[below left, pos=0.25] {event occurs} (rdy);
    \draw[observed-edge] (run) -- node[above] {release} (ext);
    \draw[observed-edge] (rdy.200) -- node[left, pos=0.55] {suspend} (rs.60);
    \draw[observed-edge] (rs.30) -- node[right, pos=0.45] {activate} (rdy.240);
    \draw[observed-edge] (blk.300) -- node[below left, pos=0.5] {suspend} (bs.135);
    \draw[observed-edge] (bs.east) -- ++(0.9,0) |- node[right, pos=0.25] {activate} (blk.east);
    \draw[observed-edge, dashed] (bs) -- node[below] {event occurs} (rs);
  \end{tikzpicture}
\caption{The seven-state model with suspension --- after Stallings, Ch.~3
\S3.2, p.~111 \cite{StallingsOS_internals_design_principles}. Top row: in
memory. Bottom row: swapped out to disk.}
\label{fig:seven-state}
\end{figure}

\section{Scheduling Queues}
\label{sec:queues}

\noindent\muted{\textit{Reading: Silberschatz et al., Ch.~3 \S3.2,
pp.~105--110.}}

A state is just a label. To act on it, the kernel keeps processes of each state
in queues, and the queues are where scheduling happens. Silberschatz et al.
describe three kinds:

\begin{definitionbox}{Scheduling queues}
The \key{job queue} holds every process in the system (in the batch-oriented
description, jobs submitted and awaiting admission). The \key{ready queue}
holds the processes in main memory that are \texttt{ready}, usually as a
linked list of PCBs. Each device has a \key{device queue} holding the
processes \texttt{waiting} for that device.
\end{definitionbox}

The implementation is simple. A queue header holds pointers to the first and
last PCB, and each PCB has a pointer field to the next PCB in the same queue.
Moving a process between queues moves a pointer, not the process.
Figure~\ref{fig:lists} shows a ready queue holding \texttt{P1}, \texttt{P4},
and \texttt{P5} and a disk queue holding \texttt{P2} and \texttt{P6}. A PCB
sits on exactly one of these queues at a time.

\begin{figure}[h]
\centering
\begin{tikzpicture}[every node/.style={font=\small},
      hd/.style={draw=primary-gold, fill=white, minimum width=2.4cm, minimum height=0.7cm, align=center, font=\scriptsize},
      pcb/.style={dag-node, minimum width=1.7cm, minimum height=0.7cm}]
    \node[hd] (rh) at (-5.0,0.9) {ready queue\\head / tail};
    \node[pcb] (a1) at (-2.2,0.9) {PCB \texttt{P1}};
    \node[pcb] (a2) at (0.2,0.9) {PCB \texttt{P4}};
    \node[pcb] (a3) at (2.6,0.9) {PCB \texttt{P5}};
    \node[hd] (dh) at (-5.0,-0.9) {disk queue\\head / tail};
    \node[pcb] (b1) at (-2.2,-0.9) {PCB \texttt{P2}};
    \node[pcb] (b2) at (0.2,-0.9) {PCB \texttt{P6}};
    \draw[observed-edge] (rh) -- (a1);
    \draw[observed-edge] (a1) -- (a2);
    \draw[observed-edge] (a2) -- (a3);
    \draw[observed-edge] (dh) -- (b1);
    \draw[observed-edge] (b1) -- (b2);
    \draw[counterfactual-edge] (rh.north) -- ++(0,0.45) -| node[above, pos=0.25, font=\scriptsize] {tail} (a3.north);
    \draw[counterfactual-edge] (dh.south) -- ++(0,-0.45) -| node[below, pos=0.25, font=\scriptsize] {tail} (b2.south);
    \node[font=\scriptsize, text=neutral] at (4.6,0.9) {next = nil};
    \node[font=\scriptsize, text=neutral] at (2.2,-0.9) {next = nil};
  \end{tikzpicture}
\caption{The ready queue and a device queue as linked lists of PCBs --- after
Silberschatz et al., Ch.~3 \S3.2, pp.~105--110.}
\label{fig:lists}
\end{figure}

\subsection{The queueing diagram}

Figure~\ref{fig:queueing} is the standard picture of how a process moves among
these queues. A new process is put in the ready queue and waits there until it
is dispatched. Once on the CPU, it leaves for exactly one of four reasons:
\begin{enumerate}
  \item it issues an \key{I/O request} and is placed on a device queue;
  \item its \key{time slice expires}, and a timer interrupt returns it to the
        ready queue;
  \item it \key{creates a child} process and may wait for the child to
        terminate;
  \item it is \key{removed by an interrupt}, or waits for an event.
\end{enumerate}
In every case except termination, the path leads back to the ready queue. In
our running example, \texttt{P2} (the editor) leaves by route~1 almost every
burst, \texttt{P1} (\texttt{gcc}) leaves mostly by route~2, and a shell running
\texttt{gcc} leaves by route~3.

\begin{figure}[h]
\centering
\begin{tikzpicture}[every node/.style={font=\scriptsize},
      bx/.style={dag-node, minimum height=0.6cm, font=\scriptsize}]
    \node[bx, minimum width=3.0cm] (rq) at (-3.2,1.6) {ready queue};
    \node[bx, minimum width=1.6cm, draw=primary-gold] (cpu) at (1.6,1.6) {CPU};
    \draw[observed-edge] (rq) -- (cpu);
    \draw[observed-edge] (cpu) -- (3.9,1.6) node[right] {\texttt{terminated}};
    \node[bx, minimum width=2.4cm] (io)  at (0.4,0.6)  {I/O queue $\to$ I/O};
    \node[bx, minimum width=2.4cm] (ts)  at (0.4,-0.2) {time slice expired};
    \node[bx, minimum width=2.4cm] (fk)  at (0.4,-1.0) {child created, runs};
    \node[bx, minimum width=2.4cm] (it)  at (0.4,-1.8) {wait for an interrupt};
    \foreach \b in {io, ts, fk, it} {
      \draw[thick, draw=jet] (3.0,1.6) |- (\b.east);
      \draw[observed-edge] (\b.west) -- (-5.4,0|-\b.west) |- (rq.west);
    }
  \end{tikzpicture}
\caption{The queueing diagram --- after Silberschatz et al., Ch.~3 \S3.2,
pp.~105--110. Every path out of the CPU, except termination, leads back to the
ready queue.}
\label{fig:queueing}
\end{figure}

\begin{example}[Queue membership over time]
\label{ex:membership}
Table~\ref{tab:membership} follows the three processes through six events on
the one-core machine.

\begin{table}[h]
\centering
\small
\renewcommand{\arraystretch}{1.2}
\begin{tabular}{@{}lllll@{}}
\toprule
\textbf{Time} & \textbf{Event} & \textbf{Running} & \textbf{Ready queue} & \textbf{Keyboard queue} \\
\midrule
$t_0$ & start & \texttt{P1} & \texttt{P2} & --- \\
$t_1$ & \texttt{P1}'s time slice expires & \texttt{P2} & \texttt{P1} & --- \\
$t_2$ & \texttt{P2} calls \texttt{read} & \texttt{P1} & --- & \texttt{P2} \\
$t_3$ & key pressed (interrupt) & \texttt{P1} & \texttt{P2} & --- \\
$t_4$ & \texttt{P1}'s time slice expires & \texttt{P2} & \texttt{P1} & --- \\
$t_5$ & \texttt{P3} admitted & \texttt{P2} & \texttt{P1}, \texttt{P3} & --- \\
\bottomrule
\end{tabular}
\caption{Queue membership for \texttt{P1}--\texttt{P3}.}
\label{tab:membership}
\end{table}

Step by step: at $t_1$ the timer preempts \texttt{P1}, which joins the ready
queue, and \texttt{P2} is dispatched. At $t_2$, \texttt{P2} blocks on
\texttt{read}, so it moves to the keyboard queue and \texttt{P1} is dispatched
again. At $t_3$ the keyboard interrupt moves \texttt{P2} to the ready queue,
where nothing is ahead of it --- but \texttt{P1} keeps the CPU until its slice
expires at $t_4$. At $t_5$ the backup job \texttt{P3} is admitted and joins the
tail of the ready queue behind \texttt{P1}. \emph{Check:} every row has exactly
one running process, and each PCB appears in exactly one place.
\end{example}

\section{Schedulers}
\label{sec:schedulers}

\noindent\muted{\textit{Reading: Silberschatz et al., Ch.~3 \S3.2,
pp.~105--110.}}

Queues hold the processes; schedulers decide who moves between them. An
operating system has up to three schedulers, and they work on very different
timescales, as Table~\ref{tab:schedulers} summarizes.

\begin{table}[h]
\centering
\small
\renewcommand{\arraystretch}{1.3}
\begin{tabular}{@{}>{\raggedright\arraybackslash}p{0.18\linewidth}>{\raggedright\arraybackslash}p{0.28\linewidth}>{\raggedright\arraybackslash}p{0.18\linewidth}>{\raggedright\arraybackslash}p{0.24\linewidth}@{}}
\toprule
\textbf{Scheduler} & \textbf{Moves a process} & \textbf{Runs} & \textbf{Controls} \\
\midrule
Long-term (job) & job pool $\to$ ready queue (\texttt{new} $\to$ \texttt{ready}) & seconds to minutes & degree of multiprogramming \\
Short-term (CPU) & ready queue $\to$ CPU (\texttt{ready} $\to$ \texttt{running}) & every few ms & who runs next \\
Medium-term & memory $\leftrightarrow$ disk (swap out / swap in) & occasionally & memory pressure, mix \\
\bottomrule
\end{tabular}
\caption{The three schedulers.}
\label{tab:schedulers}
\end{table}

\subsection{The long-term scheduler}

\begin{definitionbox}{Long-term scheduler}
The long-term (job) scheduler selects processes from the job pool on disk and
loads them into memory for execution. It controls the degree of
multiprogramming $n$, the number of processes in memory.
\end{definitionbox}

In a stable system the average rate of process admission must equal the
average rate of departure, so the long-term scheduler may need to be invoked
only when a process leaves. Because it runs rarely, it can afford a careful
choice. Batch systems have one; time-sharing systems such as UNIX and Windows
often have \emph{none}, and every new process goes straight to the short-term
scheduler.

The most important choice the long-term scheduler makes is the \key{process
mix}. An \key{I/O-bound} process spends more time doing I/O than computing and
has many short CPU bursts --- \texttt{P2}, the editor, is one. A
\key{CPU-bound} process spends most of its time computing and has few, long
bursts --- \texttt{P1}, \texttt{gcc}, is one. If every admitted process is
I/O-bound, the ready queue is almost always empty and the CPU idles; if every
process is CPU-bound, the device queues are empty and the devices idle. A good
mix keeps both busy (Silberschatz et al., Ch.~3 \S3.2, pp.~105--110).

The absence of a long-term scheduler has a visible consequence. When students
open dozens of browser tabs until a Linux laptop slows to a crawl, no scheduler
refuses them: admission control would be the long-term scheduler's job, but a
time-sharing system typically has none. Every new process goes straight to the
ready queue, and stability depends on users noticing and quitting programs.
Memory pressure is handled instead by swapping or paging --- a medium-term
mechanism --- which is why the machine slows down rather than refusing work.

\subsection{The short-term scheduler}

\begin{definitionbox}{Short-term scheduler}
The short-term (CPU) scheduler selects one process from the ready queue and
allocates the CPU to it.
\end{definitionbox}

It runs very often: at least once every time slice (typically every few tens of
milliseconds) and after every block. It must therefore be fast, because its
own running time is pure overhead.

\begin{example}[Scheduler overhead]
The short-term scheduler takes $1$\,ms to decide which process to run, and the
chosen process then runs for $10$\,ms. Each cycle is $10 + 1 = 11$\,ms long,
of which $1$\,ms is spent scheduling. The fraction of CPU time spent scheduling
is therefore
\[
  \frac{1}{10 + 1} = \frac{1}{11} \approx 9\%.
\]
\end{example}

\emph{How} the short-term scheduler chooses --- FCFS, SJF, RR, and the rest ---
is the subject of Lectures 8 and 9.

\subsection{The medium-term scheduler}

\begin{definitionbox}{Medium-term scheduler}
The medium-term scheduler removes a process from memory (swaps it out) to
reduce the degree of multiprogramming, and later reintroduces it (swaps it in)
so that it continues where it left off.
\end{definitionbox}

Swapping is used to improve the process mix or to free memory when demand
exceeds supply. Figure~\ref{fig:swap} adds it to the queueing picture.

\begin{figure}[h]
\centering
\begin{tikzpicture}[every node/.style={font=\scriptsize},
      bx/.style={dag-node, minimum height=0.7cm, font=\scriptsize}]
    \node[bx, minimum width=2.6cm] (rq) at (-3.0,0.6) {ready queue};
    \node[bx, minimum width=1.4cm, draw=primary-gold] (cpu) at (1.4,0.6) {CPU};
    \node[bx, minimum width=3.6cm, text width=3.4cm, align=center]
      (sw) at (-1.3,-1.4) {partially executed,\\swapped-out processes};
    \draw[observed-edge] (rq) -- (cpu);
    \draw[observed-edge] (cpu) -- (3.4,0.6) node[right] {end};
    \draw[observed-edge] (sw.west) -- (-5.2,-1.4) |- node[left, pos=0.25] {swap in} (rq.west);
    \draw[observed-edge] (cpu.south) |- node[right, pos=0.3] {swap out} (sw.east);
  \end{tikzpicture}
\caption{The medium-term scheduler adds a swap-out/swap-in loop --- after
Silberschatz et al., Ch.~3 \S3.2, pp.~105--110.}
\label{fig:swap}
\end{figure}

\subsection{Schedulers on the state diagram}

Putting the two diagrams together shows exactly which transitions each scheduler
drives. The long-term scheduler drives \texttt{new} $\to$ \texttt{ready}
(admission). The short-term scheduler drives \texttt{ready} $\to$
\texttt{running} (dispatch) and, with preemption, also decides
\texttt{running} $\to$ \texttt{ready}. The medium-term scheduler drives
\texttt{ready} $\leftrightarrow$ \texttt{ready/suspend} and \texttt{blocked}
$\leftrightarrow$ \texttt{blocked/suspend}. Two transitions are \emph{not} a
scheduler's decision: \texttt{running} $\to$ \texttt{waiting} is caused by the
process's own request, and \texttt{waiting} $\to$ \texttt{ready} is caused by
the completion interrupt.

\section{The Context Switch}
\label{sec:cs}

\noindent\muted{\textit{Reading: Silberschatz et al., Ch.~3 \S3.2, p.~110;
Stallings, Ch.~3 \S3.4, p.~135; Bovet \& Cesati, Ch.~3, pp.~79--130.}}

The short-term scheduler picks the next process. Handing it the CPU means
saving one process's context and restoring another's.

\begin{definitionbox}{Context and context switch}
The context of a process is the CPU and kernel state needed to resume it:
register values, \texttt{PC}, process state, and memory-management
information, all represented in its PCB. A context switch saves the context of
the current process into its PCB and restores the saved context of the process
scheduled to run next.
\end{definitionbox}

A context switch has two halves, a \key{state save} and a \key{state restore}.
It is triggered by an interrupt, a trap, or a system call that blocks.
Figure~\ref{fig:timeline} shows two processes alternating: an interrupt moves
the CPU from \texttt{P0} to the kernel, which saves PCB$_0$ and reloads
PCB$_1$; later a system call from \texttt{P1} triggers the reverse. While the
kernel switches, neither process executes.

\begin{figure}[h]
\centering
\begin{tikzpicture}[every node/.style={font=\scriptsize}]
    \node[anchor=east] at (-5.2,1.2) {\texttt{P0}};
    \node[anchor=east] at (-5.2,0) {OS};
    \node[anchor=east] at (-5.2,-1.2) {\texttt{P1}};
    % P0 executing / idle
    \draw[line width=3pt, draw=primary-blue] (-5,1.2) -- (-2.6,1.2);
    \node[above] at (-3.8,1.3) {executing};
    \draw[thick, dashed, draw=neutral] (-2.6,1.2) -- (4.6,1.2);
    \node[above] at (1.0,1.3) {idle};
    \draw[line width=3pt, draw=primary-blue] (4.6,1.2) -- (6.0,1.2);
    \node[above] at (5.3,1.3) {executing};
    % OS work
    \node[dag-node, minimum width=2.0cm, minimum height=0.7cm, font=\scriptsize,
      draw=primary-gold, align=center] (s1) at (-1.4,0) {save PCB$_0$\\reload PCB$_1$};
    \node[dag-node, minimum width=2.0cm, minimum height=0.7cm, font=\scriptsize,
      draw=primary-gold, align=center] (s2) at (3.4,0) {save PCB$_1$\\reload PCB$_0$};
    % P1 idle / executing / idle
    \draw[thick, dashed, draw=neutral] (-5,-1.2) -- (-0.2,-1.2);
    \node[below] at (-2.6,-1.3) {idle};
    \draw[line width=3pt, draw=primary-blue] (-0.2,-1.2) -- (2.2,-1.2);
    \node[below] at (1.0,-1.3) {executing};
    \draw[thick, dashed, draw=neutral] (2.2,-1.2) -- (6.0,-1.2);
    \node[below] at (4.6,-1.3) {idle};
    % trigger arrows
    \draw[observed-edge] (-2.6,1.05) -- node[left, pos=0.45] {interrupt} (s1.north west);
    \draw[observed-edge] (s1.south east) -- (-0.2,-1.05);
    \draw[observed-edge] (2.2,-1.05) -- node[left, pos=0.4] {system call} (s2.south west);
    \draw[observed-edge] (s2.north east) -- (4.6,1.05);
  \end{tikzpicture}
\caption{A context switch on a timeline --- after Silberschatz et al., Ch.~3,
pp.~103--110. The gold boxes are pure overhead.}
\label{fig:timeline}
\end{figure}

\subsection{The steps of a context switch}

Stallings (Ch.~3 \S3.4, p.~135) lists the full sequence of process-switch steps;
Bovet \& Cesati (Ch.~3, pp.~79--130 \cite{BovetCesati_understanding_linux_kernel})
describe Linux's implementation. In outline, switching from \texttt{P0} to
\texttt{P1}:
\begin{enumerate}
  \item An interrupt or trap enters the kernel; the hardware saves the user
        \texttt{PC} and flags.
  \item The kernel saves the remaining registers of \texttt{P0} into PCB$_0$.
  \item It updates PCB$_0$'s state (\texttt{ready} or \texttt{waiting}) and
        links PCB$_0$ onto the right queue.
  \item The short-term scheduler selects \texttt{P1}.
  \item The kernel sets PCB$_1$'s state to \texttt{running} and switches
        memory-management state (page-table base, TLB as needed).
  \item It restores \texttt{P1}'s registers from PCB$_1$ and returns to user
        mode at \texttt{P1}'s saved \texttt{PC}.
\end{enumerate}

\subsection{The dispatcher}

\begin{definitionbox}{Dispatcher}
The dispatcher is the kernel module that gives control of the CPU to the
process selected by the short-term scheduler. The time it takes to stop one
process and start another is the dispatch latency.
\end{definitionbox}

The dispatcher switches context, switches to user mode, and jumps to the proper
location in the user program. The scheduler is the \emph{policy} (which
process?); the dispatcher is the \emph{mechanism} (make it run). Every dispatch
involves a context switch. The dispatcher is previewed here and developed in
Lecture~8 (Silberschatz et al., Ch.~5 \S5.1, from p.~183).

\subsection{A mode switch is not a context switch}

A \key{mode switch} goes user $\to$ kernel $\to$ user for the \emph{same}
process --- for example, a \texttt{getpid} system call. A few registers are
saved and restored, and the PCB state does not change. A \key{context (process)
switch} moves the CPU to a \emph{different} process: a full register save, a
PCB state update, a queue move, and a memory-map switch. Every context switch
includes mode switches, but most mode switches cause no context switch; an
interrupt can be handled and return to the interrupted process with no process
switch at all (Stallings, Ch.~3 \S3.4, p.~135).

The test is simple: \emph{does the caller keep the CPU when the kernel returns?}
A process that calls \texttt{getpid} does --- the kernel answers and returns,
a mode switch only. A process that calls \texttt{read} on an empty pipe does
not: it must block (\texttt{running} $\to$ \texttt{waiting}), so the kernel
switches to another process.

\subsection{The cost of switching}

A context switch is pure overhead: the system does no useful work while
switching. Its speed depends on the hardware --- memory speed, the number of
registers to copy, and special instructions such as a single instruction to
load or store all registers. A typical cost is a few microseconds, plus an
indirect cost from cold caches and TLB. Some processors provide multiple
register sets, so a switch simply changes a pointer to the current set. The
general rule (Silberschatz et al., Ch.~3 \S3.2, p.~110): the more complex the
operating system and its memory management, the more context must be saved and
the more each switch costs.

If a switch costs $s$ and each process runs for a time slice $q$ before being
switched out, each cycle lasts $q + s$, of which $s$ is overhead:
\[
  \text{overhead fraction} = \frac{s}{q + s}.
\]

\begin{example}[Overhead of switching]
\label{ex:overhead}
Let $s = 5\ \mu\text{s}$. The number of switches per second is approximately
$1/(q+s)$.
\begin{itemize}
  \item $q = 10$\,ms $= 10000\ \mu$s: $1/(10005\ \mu\text{s}) \approx 99.95$
        switches/s; overhead $5/10005 \approx 0.05\%$.
  \item $q = 1$\,ms $= 1000\ \mu$s: $\approx 995$ switches/s; overhead
        $5/1005 \approx 0.50\%$.
  \item $q = 100\ \mu$s: $\approx 9524$ switches/s; overhead
        $5/105 \approx 4.76\%$.
  \item $q = 10\ \mu$s: $\approx 66667$ switches/s; overhead
        $5/15 \approx 33.3\%$.
\end{itemize}
\emph{Conclusion:} shrink the slice and the CPU spends more of its life
switching than working.
\end{example}

\begin{example}[\texttt{P1} and \texttt{P2} on one core]
\label{ex:p1p2}
\begin{enumerate}
  \item $t = 0$: \texttt{P1} (\texttt{gcc}) is running; \texttt{P2} is waiting
        on the keyboard.
  \item $t = 3$\,ms: a key is pressed. The interrupt moves \texttt{P2} to
        \texttt{ready}. This is a mode switch only: \texttt{P1} resumes.
  \item $t = 10$\,ms: \texttt{P1}'s slice expires. Context switch \#1: save
        PCB$_1$, \texttt{P1} $\to$ \texttt{ready}; restore PCB$_2$,
        \texttt{P2} $\to$ \texttt{running}.
  \item $t = 10.2$\,ms: \texttt{P2} echoes the key and calls \texttt{read}.
        Context switch \#2: \texttt{P2} $\to$ \texttt{waiting}; \texttt{P1}
        $\to$ \texttt{running}.
\end{enumerate}
\emph{Count:} two context switches and three kernel entries (interrupt, timer,
system call) in $10.2$\,ms. With $s = 5\ \mu$s the switching cost is
$2 \times 5 = 10\ \mu\text{s}$, and $10/10200 \approx 0.1\%$ of the interval.
\end{example}

\section{Putting It Together}

States, queues, schedulers, and the context switch form one machine. Following
the backup job \texttt{P3} through its whole life uses every part of it.

\begin{example}[\texttt{P3}'s life, end to end]
\begin{enumerate}
  \item Submitted to the batch queue: \texttt{new}, in the job queue.
  \item The long-term scheduler admits it at night: \texttt{new} $\to$
        \texttt{ready}, PCB on the ready queue.
  \item The short-term scheduler picks it and the dispatcher context-switches
        to it: \texttt{ready} $\to$ \texttt{running}.
  \item It writes to the backup disk: \texttt{running} $\to$ \texttt{waiting},
        on the disk queue; the kernel context-switches away.
  \item Memory is short, so the medium-term scheduler swaps it out:
        \texttt{blocked} $\to$ \texttt{blocked/suspend}.
  \item The write completes: \texttt{blocked/suspend} $\to$
        \texttt{ready/suspend}; later it is swapped in: $\to$ \texttt{ready}.
  \item It is dispatched again, finishes, and calls \texttt{exit}:
        \texttt{running} $\to$ \texttt{terminated}.
\end{enumerate}
All three schedulers and both diagrams appear in one life.
\end{example}

\subsection{Common misconceptions}

\begin{itemize}
  \item \bad{``A woken process starts running immediately.''} It goes
        \texttt{waiting} $\to$ \texttt{ready} and queues.
  \item \bad{``Every system call is a context switch.''} Most are mode
        switches; only blocking or preemption switches process.
  \item \bad{``Every OS has a long-term scheduler.''} Time-sharing systems
        often have none.
  \item \bad{``The ready queue holds waiting processes.''} Waiting processes are
        on device queues.
\end{itemize}

\subsection{Summary}

\begin{itemize}
  \item A process is in one of five states --- \texttt{new}, \texttt{ready},
        \texttt{running}, \texttt{waiting}, \texttt{terminated} --- and six
        transitions are legal. Suspension adds two more states.
  \item The kernel keeps PCBs on a job queue, a ready queue, and device queues;
        the queueing diagram shows every path.
  \item The long-term scheduler sets the degree of multiprogramming; the
        short-term scheduler picks the next process every few milliseconds; the
        medium-term scheduler swaps.
  \item A context switch saves one PCB and restores another; it is pure
        overhead, a fraction $s/(q+s)$ of the CPU.
  \item A mode switch is not a context switch.
\end{itemize}

\paragraph{Next.} This chapter dealt with processes that already exist. The
next lecture asks how processes come to exist and how they end ---
\texttt{fork}, \texttt{exec}, \texttt{wait}, \texttt{exit}, and the process
tree. Reading: Silberschatz et al., Ch.~3 \S3.3
\cite{SilberschatzOSC_operating_system_concepts}; Stallings, Ch.~3 \S3.4
\cite{StallingsOS_internals_design_principles}; Bovet \& Cesati, Ch.~3
\cite{BovetCesati_understanding_linux_kernel}.

For this chapter, the page-cited sources are Silberschatz et al., 8th ed., Ch.~3
\S3.1 (process states, p.~103) and \S3.2 (process scheduling: queues,
schedulers, context switch, pp.~105--110); Stallings, 6th ed., Ch.~3 \S3.2
(two-, five-, and seven-state models, p.~111) and \S3.4 (process control, mode
and process switching, p.~135); and Bovet \& Cesati, 3rd ed., Ch.~3 (processes
and the process switch, pp.~79--130). All are page-verified through each book's
\texttt{index.md} in \texttt{master\_supporting\_docs/CS502/supporting\_books/}.

\section*{Self-Check}
\addcontentsline{toc}{section}{Self-Check}

\subsection*{Exercises}

\begin{enumerate}
  \item For each pair, state whether the transition is legal in the five-state
        model and, if it is, name the event and the agent: (a) \texttt{ready}
        $\to$ \texttt{running}; (b) \texttt{waiting} $\to$ \texttt{running};
        (c) \texttt{running} $\to$ \texttt{waiting}; (d) \texttt{ready} $\to$
        \texttt{waiting}.
  \item An 8-core server has 120 processes. At one instant 8 are running, 17
        are ready, and 2 are new. How many are waiting, assuming none is
        terminated? What is the largest possible number of running processes?
  \item A short-term scheduler takes $0.5$\,ms per decision and gives each
        process a $4.5$\,ms slice. What fraction of the CPU is spent
        scheduling?
  \item A context switch costs $s = 2\ \mu$s. Find the overhead fraction for
        $q = 1$\,ms and $q = 20\ \mu$s, and the slice $q$ at which overhead
        reaches exactly $10\%$.
  \item A process calls \texttt{getpid}, then \texttt{read} on a disk file
        whose block is not cached, then \texttt{exit}. For each call say
        whether a context switch must follow, and give the state transitions.
\end{enumerate}

\subsection*{Solutions}

\begingroup\small
\begin{enumerate}
  \item (a) Legal: the short-term scheduler selects it and the dispatcher loads
        its context. (b) Illegal: a woken process goes \texttt{waiting}
        $\to$ \texttt{ready} and must be dispatched. (c) Legal: the process
        itself issues an I/O request or \texttt{wait} (a system call).
        (d) Illegal: a process that is not running cannot issue the request it
        would wait on.
  \item Waiting $= 120 - 8 - 17 - 2 = 93$. With 8 cores, at most 8 processes
        can be running at once.
  \item Each cycle lasts $4.5 + 0.5 = 5$\,ms, of which $0.5$\,ms is scheduling:
        $0.5/5 = 10\%$.
  \item $q = 1000\ \mu$s: $2/1002 \approx 0.20\%$. $q = 20\ \mu$s:
        $2/22 \approx 9.09\%$. Setting $s/(q+s) = 0.1$ gives $q + s = 10s$, so
        $q = 9s = 18\ \mu$s.
  \item \texttt{getpid}: mode switch only; the process stays \texttt{running}.
        \texttt{read} of an uncached block: the process must wait for the disk,
        so \texttt{running} $\to$ \texttt{waiting} and the kernel context-switches
        to another process; on completion \texttt{waiting} $\to$ \texttt{ready},
        and later \texttt{ready} $\to$ \texttt{running} when dispatched.
        \texttt{exit}: \texttt{running} $\to$ \texttt{terminated}; the process
        never runs again, so the kernel must switch to another process.
\end{enumerate}
\endgroup

\begin{thebibliography}{9}
\bibitem{SilberschatzOSC_operating_system_concepts}
A.~Silberschatz, P.~B.~Galvin, and G.~Gagne, \emph{Operating System Concepts},
8th ed., John Wiley and Sons --- Ch.~3 \S3.1 (process states, p.~103) and
\S3.2 (process scheduling, pp.~105--110). \emph{(Page-cited; index built
2026-10-14.)}
\bibitem{StallingsOS_internals_design_principles}
W.~Stallings, \emph{Operating Systems: Internals and Design Principles}, 6th
ed., Prentice Hall --- Ch.~3 \S3.2 (process states, p.~111) and \S3.4
(process control, p.~135). \emph{(Page-cited; index built 2026-10-14.)}
\bibitem{BovetCesati_understanding_linux_kernel}
D.~P.~Bovet and M.~Cesati, \emph{Understanding the Linux Kernel}, 3rd ed.,
O'Reilly Media, 2006 --- Ch.~3 (processes and the process switch,
pp.~79--130). \emph{(Page-cited; index built 2026-10-14.)}
\end{thebibliography}

\end{document}
